% P10-4c, 14.09.2026: Ledger-Herleitung und Nachweisanschlüsse präzisiert.
% M09, 11.09.2026: Leserführung und MCR konzentriert; neue
% HERKUNFT-M09-Kommentare unten sind die aktuelle Schreibgrundlage.
% Die folgende Versionshistorie beschreibt frühere Fassungen.
% ============================================================
% §4.5 Kontrollierte semantische Transformation — v08 (Claude)
% v08 (01.09., Kollegen-Freeze-Runde — alle 4 Punkte berechtigt):
% ① Designproblem ohne „gemeinsamen Verarbeitungsschritt" (Konsistenz
% mit der entschärften Ausgangslage) ② E45-02: strukturelle Cluster→
% Arbeitspaket-Regel IM TEXT disambiguiert (≠ Quellenabdeckung aus 4.4);
% dabei Kommentar-Verschlucker gefixt ③ Grenzen je Prüfform getrennt
% (modellgestützt übersieht · Repair führt ein · deterministisch sichert
% NUR Formalisiertes) ④ Aufwand differenziert (Modellaufrufe nur bei
% modellgestützten Schritten + Orchestrierungsaufwand). A4 → weiterhin
% Gesamtcheck-Merker (Kollege bestätigt Vorgehen).
% v07 (01.09., Freeze-Kandidaten-Runde): ① E45-03-Einzelfall PRÄZISIERT
% — B42 nachgelesen: der belegte Fall war KEINE Auslassung, sondern eine
% als nicht quellgedeckt EINGESTUFTE Technologiefestlegung (Kollegen-
% Einwand berechtigt und durch die Note verschärft; Anhang-Zeile
% mitgezogen). ② Maker-Eingabeseite ergänzt (Rolle+Eingabe+Ausgabe).
% ③ „Agentenknoten bleiben agentisch" → „Einbettung hebt den Agenten-
% charakter nicht auf" (2.2-Anschluss).
% v06 (01.09., Kollegen-Runde zu v05 — KRITISCH geprüft): übernommen
% ① Rückkante Repair→Checker + Terminal-Stopp (R-33-belegt) ② MCR-
% Schleife vs. Human-Gate/Apply getrennt ③ Checker = kriteriums-
% abhängige Verantwortung ④ E45-03-Übergangssatz ⑤ Ausgangslage
% entschärft (Same-Call-Selbstprüfung für diese Ära nicht belegbar).
% NICHT übernommen: A4-Umformulierung (Blaupausen-Kanon + 4.8-Matrix-
% Kopplung → MERKER für den Autor-Gesamt-Check über A1–A8; Kollege
% selbst: „später gegen A1–A8 insgesamt prüfen"); „strukturierter"
% statt „wörtlicher" Befund (wörtlich ist die belegte Mechanik).
% v05 (01.09., Erstleser-Runde auf Autor-Wunsch): ① E44-01-Satz
% erzählt statt komprimiert (BEIDE Prüfarten modellgetragen — freie
% Fehlersuche vs. feste Checkliste; „modellabhängig/verrauscht/
% Befundmengen"-Jargon raus). ② Muster-Stationen = eigenständige
% Workflow-Knoten (2.4.1-Anschluss; Autor-Punkt „Schritte sind
% eigentlich Workflows"). ③ Akteurs-Markierung im Muster komplett:
% Maker/Repair agentisch · Checker det./modellgestützt je Kriterium ·
% Autorisierung Mensch · Anwendung deterministisch; Mega-Satz in
% vier Sätze geteilt.
% v04 (01.09., Konkretisierungs-Runde): M7 Repair-Vignette (70×
% UNCOVERED_CORE = struktureller Cluster→PBI-Schnitt-Check, Pass in
% Attempt 2; Grenz-Satz „Rückkopplungspfad, nicht Korrektheit");
% E45-02-ID in die Vignette verschoben.
% v02 (01.09., Kollegen-Review): E45-03-Neufassung (Variabilität +
% dokumentierter Einzelfund statt Wertung „inhaltlich schärfer");
% „mehrere der betrachteten Fehlerklassen"; E44-01-Wiederverwendung.
% v03 (01.09., „dosierte Agentik"-Runde): Gestaltungsprinzip-Absatz
% nach A4 (semantisch offen = agentisch · formal ausdrückbar =
% deterministisch · Workflow verbindet; Agentenknoten bleiben agentisch).
%
% Grundlage: Kapitel-4-Blaupause §10 (Absatzfolge 1–7, A4) ·
% genese-dossier Z6/P4 (AUDITIERT) · Beweiskette Säule 2.
%
% HERKUNFTS-ENTSCHEID (Blaupause §10 "Benötigte Evidenz"): Die
% 4.8-Matrix deklariert A4 als "formative Beobachtung beziehungsweise
% Designrationale" — der Text macht diese MISCHFORM explizit:
%   · FORMATIV belegt ist die Prüf-Seite: D1-Spike (committet) —
%     freie modellgetragene Prüfung modellabhängig + verrauscht,
%     feste Prüfliste über Modelle hinweg gleichartig (S2-Richtung,
%     KEINE Zahlen; "Severity 6→2" bleibt draußen).
%   · RATIONALE ist die Rollen-Seite (Autorisierung/Anwendung):
%     folgt aus den Kontrollebenen 2.3.3/2.3.4, nicht aus einem
%     Fehlschlag — im Text ausdrücklich so gekennzeichnet (kein
%     erfundener Pivot!).
%
% Audit-Anker (Anhang-Tabelle; committet ✓):
%   D1-direct-vs-topic (Prüf-Modellabhängigkeit) · dokumentierter
%   Loop-Lauf 074419 (Prüf-Feedback → Reparatur → Bestehen; Chronik)
%   als Beleg der etablierten Musterform.
%
% Regeln: "Maker–Checker–Repair" hier als Musterbegriff EINGEFÜHRT
% (Blaupause §10 sieht das vor); keine Instanzen-Aufzählung (Kap. 5);
% keine Mess-Verben; Aufwands-Frage explizit nach Kap. 7 verwiesen
% (Vergleich mit/ohne Prüf- und Reparaturschritte). Budget ~1,2 S.
% ============================================================

\section{Kontrollierte semantische Transformation}
\label{sec:kontrollierte-semantische-transformation}

\emph{Ausgangslage.} Die Quellenstrukturierung aus
Abschnitt~\ref{sec:evidenzgebundene-wissensgewinnung} machte
Zwischenergebnisse überprüfbar, entschied aber noch nicht über ihre
Übernahme. Beim Bilden quellengebundener Aussagen und Ableiten
fachlicher Vorschläge konnten sowohl inhaltliche Fehler als auch
ungültige Referenzen oder fehlende Pflichtangaben entstehen. Dafür
waren unterschiedliche Prüfmittel erforderlich.

\emph{Formativer Befund und Designproblem.} Der Themenreview aus
Abschnitt~\ref{sec:evidenzgebundene-wissensgewinnung} zeigte den
Nutzen begrenzter Prüfaufträge im untersuchten Fall. Daraus ergab
sich die weiterführende Gestaltungsfrage, welche Urteile Modelle
übernehmen sollten und welche Bedingungen sich formal prüfen ließen.
% HERKUNFT M09: D1-direct-vs-topic/evidence.md und acht R1-JSONs
% gegengeprüft. Missing-Zahlen stark 1/4/5/6, mini 1/2/3/5;
% beide Litmus-Themen unter MissingTopic in open-questions, nicht
% identische Gesamtbefunde oder identische Artefaktzuordnung.

Ein weiterer Vergleich zeigte zugleich den Nutzen offener
Beurteilung (E45-03). Bei drei Prüfungen desselben Risikoartefakts
variierte der frei prüfende Agent in Anzahl, Kategorien und
Granularität seiner Befunde. Er benannte aber auch eine
Technologiefestlegung als nicht durch die Quelle gedeckt, die die
zerlegte Prüfstrecke nicht angezeigt hatte. Dies war die Einstufung
des Prüfers, kein unabhängig bestätigter Fehler. Der Fall begründete
somit keine generelle Überlegenheit einer Prüfungsform.
% HERKUNFT M09: iteration-notes-Phase2B-Commited.md B42;
% drei risks.review.md aus runsArchive/reviewpilot, 8/6/7 Befunde.

Aus diesen Beobachtungen wurde eine Aufgabenverteilung abgeleitet:
Modelle sollten Bedeutung interpretieren und fachliche Vorschläge
beurteilen; bekannte Struktur-, Referenz- und Pflichtfeldregeln
sollten deterministisch geprüft werden. Ein weiterer Modellaufruf
war für diese formal ausdrückbaren Regeln nicht erforderlich.
Entscheidend war damit nicht allein die Wahl eines stärkeren
Modells, sondern auch, welche Verantwortung ihm übertragen wurde.

\emph{Entwurfsentscheidung.} Erzeugung, Prüfung und eine an
Prüfbefunde gebundene Reparatur wurden als getrennte Workflow-Schritte
realisiert. Dieses Maker--Checker--Repair-Muster ermöglicht eine
begrenzte Überarbeitung mit anschließender erneuter Prüfung;
Abschnitt~\ref{sec:saeule-semantische-transformation} erklärt den
Ablauf. Ein späterer Lauf dokumentiert seine Ausführung an einer
strukturellen Zuordnungsregel für Arbeitspakete (E45-02), ohne damit
die fachliche Reparaturqualität nachzuweisen. Die zusätzliche
Trennung von menschlicher Autorisierung und Zustandsanwendung
überträgt die Kontrollebenen aus den
Abschnitten~\ref{sec:human-in-the-loop-oversight} und
\ref{sec:governance-autorisierung-kontrollen} auf das System; sie
wird nicht nachträglich als Reaktion auf einen beobachteten
Fehlschlag ausgegeben. Daraus ergibt sich die vierte
Designanforderung:

\textbf{A4:} Semantische Erzeugung beziehungsweise Transformation,
Prüfung, Reparatur, fachliche Autorisierung und Zustandsmutation sind
als getrennte Verantwortlichkeiten zu realisieren.

Der explizite Workflow verbindet diese Aufgaben und bestimmt
Reihenfolge, Wiederholung und Übergänge. Agentische Entscheidungen
bleiben dort möglich, wo Interpretation oder Auswahl erforderlich
sind; die Ablaufsteuerung muss dafür nicht ebenfalls einem Modell
überlassen werden (Abschnitt~\ref{sec:orchestrierungsformen}).

\emph{Einordnung und Grenzen.} Die Trennung macht Prüfbefunde und
Übergänge nachvollziehbar. Sie schließt weder übersehene inhaltliche
Fehler noch neue Fehler durch Reparaturen aus. Deterministische
Prüfungen erfassen nur die jeweils formalisierten Bedingungen.
Kapitel~\ref{chap:evaluation} untersucht ausgewählte Änderungen der
Inhaltsabdeckung, definierte Kontrollfälle und den Tokenaufwand der
bewerteten Konfigurationen. Daraus folgt weder ein isolierter
Wirksamkeitsnachweis jeder Prüf- oder Reparaturstufe noch ein
allgemeines Verhältnis von Aufwand und Nutzen.
